\documentclass[a4paper,12pt]{article}
\usepackage{iftex}
\ifPDFTeX
  \usepackage[T2A]{fontenc}
  \usepackage[utf8]{inputenc}
  \usepackage{helvet}
  \renewcommand{\familydefault}{\sfdefault}
\else
  \usepackage{fontspec}
  \setmainfont{Liberation Sans}
  \setsansfont{Liberation Sans}
\fi
\usepackage[russian]{babel}
\usepackage{amsmath,amssymb,mathtools}
\usepackage{icomma}
\usepackage[a4paper,left=21mm,right=21mm,top=21mm,bottom=21mm,headheight=25pt]{geometry}
\usepackage{microtype}
\usepackage{tikz}
\usetikzlibrary{arrows.meta,positioning,shapes.geometric,calc}
\usepackage{pgfplots}
\pgfplotsset{compat=1.18}
\usepackage{booktabs,tabularx}
\usepackage{enumitem}
\usepackage{fancyhdr}
\usepackage{setspace}
\usepackage{xcolor}
\usepackage{hyperref}
\definecolor{accent}{HTML}{087C82}
\definecolor{darkaccent}{HTML}{174F59}
\definecolor{textgray}{HTML}{373D42}
\definecolor{softgray}{HTML}{68747A}
\definecolor{lightline}{HTML}{CCD9DB}
\hypersetup{colorlinks=true,linkcolor=darkaccent,urlcolor=darkaccent,pdfauthor={Лёвин Артём Александрович},pdftitle={Второй закон Ньютона. Чек-лист для Володи Хачатуряна}}
\onehalfspacing
\setlength{\parindent}{0pt}
\setlength{\parskip}{5pt}
\setlist[enumerate]{leftmargin=*,itemsep=7pt,topsep=5pt}
\setlist[itemize]{leftmargin=*,itemsep=1pt,topsep=2pt}
\pagestyle{fancy}
\fancyhf{}
\fancyhead[L]{\small\color{softgray} ФИЗИКА · ОГЭ}
\fancyhead[R]{\small\color{softgray} 16.10.26}
\fancyfoot[C]{\small\color{softgray}\thepage}
\renewcommand{\headrulewidth}{0.3pt}
\renewcommand{\headrule}{\hbox to\headwidth{\color{lightline}\leaders\hrule height \headrulewidth\hfill}}
\newcommand{\heading}[1]{\vspace{2pt}{\large\bfseries\color{darkaccent}#1}\par\vspace{3pt}}
\newcommand{\subheading}[1]{\vspace{1pt}{\bfseries\color{darkaccent}#1}\par}
\newcommand{\focus}[1]{\textcolor{darkaccent}{\textbf{#1}}}
\newcommand{\siunit}[1]{\,\text{#1}}
\begin{document}
\thispagestyle{empty}
\begin{center}
\vspace*{40mm}
{\small\color{accent}\bfseries ФИЗИКА · ПОДГОТОВКА К ОГЭ}\par
\vspace{17mm}
{\Huge\bfseries\color{darkaccent} Чек-лист}\par
\vspace{8mm}
{\LARGE\bfseries Второй закон Ньютона}\par
\vspace{3mm}
{\Large Равнодействующая сила и ускорение}\par
\vspace{10mm}
{\normalsize\color{textgray} Повторение кинематики и движения по окружности}\par
\vspace{17mm}
{\large 16 октября 2026 года}\par
\vfill
{\normalsize\color{textgray} Лёвин Артём Александрович}\par
{\normalsize\color{textgray} эксклюзивно для Володи Хачатуряна}\par
\vspace{24mm}
\begin{tikzpicture}[x=1cm,y=1cm]
\draw[accent!65,line width=1.2pt] (-5,0) .. controls (-2.8,0.9) and (-1.2,-0.9) .. (1,0) .. controls (2.9,0.75) and (4.1,0.45) .. (5,0.1);
\end{tikzpicture}
\end{center}
\clearpage

{\huge\bfseries\color{darkaccent} Основной чек-лист}\par
\vspace{5pt}
\heading{1. Равнодействующая и второй закон Ньютона}
\begin{itemize}
\item[$\square$] \focus{Равнодействующая} --- векторная сумма \emph{всех} сил, действующих на тело:
\[\vec F_{\text{рез}}=\vec F_1+\vec F_2+\ldots\]
\item[$\square$] В инерциальной системе отсчёта второй закон Ньютона:
\[\boxed{\vec F_{\text{рез}}=m\vec a}\qquad \boxed{F_{\text{рез},x}=ma_x}.\]
Здесь $m$ --- масса (кг), $a$ --- ускорение ($\text{м}/\text{с}^2$), $F$ --- сила (Н).
\item[$\square$] Если $\vec a=\vec0$, то $\vec F_{\text{рез}}=\vec0$. Тело может при этом \emph{покоиться или двигаться равномерно и прямолинейно}. Отдельные силы вполне могут быть ненулевыми и уравновешивать друг друга.
\item[$\square$] Для модулей равнодействующей и ускорения: $F_{\text{рез}}=ma$. Если нужна проекция на ось, используем знаки: $F_{\text{рез},x}=ma_x$.
\end{itemize}
\subheading{Простейший пример}
Тележка массой $3\siunit{кг}$ движется с равнодействующей силой $6\siunit{Н}$. Тогда
\[a=\frac{F_{\text{рез}}}{m}=\frac{6}{3}=2\siunit{м/с\textsuperscript{2}}.\]

\heading{2. Как восстановить ускорение}
Для \textbf{постоянного ускорения} вдоль выбранной оси:
\[
\boxed{a_x=\frac{v_x-v_{0x}}{t}}\qquad
\boxed{a_x=\frac{v_{x2}-v_{x1}}{t_2-t_1}}.
\]
Вторая формула удобна для прямолинейного участка графика $v_x(t)$: ускорение равно его \emph{наклону}. Выбирайте две точно читаемые точки; для прямой подойдут любые две разные точки.
\[
\boxed{\Delta x=v_{0x}t+\frac{a_xt^2}{2}}\qquad
\boxed{v_x=v_{0x}+a_xt}.
\]
\[
\boxed{\Delta x=\frac{v_{0x}+v_x}{2}\,t}\qquad
\boxed{v_x^2-v_{0x}^2=2a_x\Delta x}.
\]
Здесь $\Delta x=x-x_0$ --- \emph{перемещение в проекции}. Если тело движется только вдоль положительного направления оси, то $s=\Delta x$ (пройденный путь). В формуле с делением на $a_x$ требуется $a_x\ne0$.
\clearpage

{\huge\bfseries\color{darkaccent} Методы и примеры}\par
\vspace{5pt}
\heading{3. Сравнение двух состояний}
Вводим индексы 1 и 2. При $m_1,m_2>0$ и сравнении модулей равнодействующих:
\[
\boxed{\frac{a_2}{a_1}=\frac{F_2}{F_1}\cdot\frac{m_1}{m_2}}\qquad
\boxed{\frac{F_2}{F_1}=\frac{m_2}{m_1}\cdot\frac{a_2}{a_1}}.
\]
Формулы отношений используют ненулевые величины в знаменателях.

\subheading{Пример: что меняется при удвоении силы и массы?}
Если $F_2=2F_1$ и $m_2=2m_1$, то $a_2/a_1=2/2=1$: ускорение прежнее.
Если же $F_2=2F_1$ и $m_2=m_1/2$, то $a_2/a_1=2/(1/2)=4$: ускорение возрастёт вчетверо. \focus{Уменьшение массы и её увеличение дают разные ответы.}

\subheading{Пример: как изменить силу?}
Массу уменьшили в 2 раза, ускорение --- в 4 раза:
\[\frac{F_2}{F_1}=\frac12\cdot\frac14=\frac18.\]
Следовательно, \textbf{силу нужно уменьшить в 8 раз}. Ответ к вопросу «во сколько раз?» записываем без единиц измерения.

\heading{4. Сила по графику скорости}
Автомобиль массой $2000\siunit{кг}$ движется с постоянным ускорением. График $v_x(t)$ проходит через точки $(4;6)$ и $(20;10)$ (время --- в секундах, скорость --- в м/с).
\begin{center}
\begin{tikzpicture}
\begin{axis}[
width=11.7cm,height=5cm,xmin=0,xmax=22,ymin=4,ymax=11,
axis lines=left,xlabel={$t$, с},ylabel={$v_x$, м/с},
grid=major,xtick={0,4,8,12,16,20},ytick={4,6,8,10},
tick label style={font=\small},label style={font=\small},
axis line style={draw=textgray},grid style={lightline!70},clip=false]
\addplot[accent,very thick,domain=0:21,samples=50]{5+x/4};
\addplot[only marks,mark=*,mark size=2.5pt,accent] coordinates {(4,6) (20,10)};
\node[above left,font=\small,fill=white,inner sep=1pt] at (axis cs:4,6) {$(4;6)$};
\node[above left,font=\small,fill=white,inner sep=1pt] at (axis cs:20,10) {$(20;10)$};
\end{axis}
\end{tikzpicture}
\end{center}
\vspace{-4pt}
\[a_x=\frac{10-6}{20-4}=0{,}25\siunit{м/с\textsuperscript{2}},\qquad
F_{\text{рез},x}=2000\cdot0{,}25=500\siunit{Н}.\]
\clearpage

{\huge\bfseries\color{darkaccent} Формулы по данным задачи}\par
\vspace{6pt}
\heading{5. Сила, если известны путь и время}
Брусок массой $1{,}10\siunit{кг}$ спускается из состояния покоя по прямой наклонной плоскости длиной $15\siunit{см}$ за $0{,}26\siunit{с}$. Считаем ускорение постоянным и направленным вдоль движения.
\[
s=15\siunit{см}=0{,}15\siunit{м},\qquad v_0=0.
\]
\[
s=\frac{at^2}{2}\quad\Longrightarrow\quad
 a=\frac{2s}{t^2}=\frac{2\cdot0{,}15}{0{,}26^2}\approx4{,}44\siunit{м/с\textsuperscript{2}}.
\]
\[
F_{\text{рез}}=ma\approx1{,}10\cdot4{,}44\approx\boxed{4{,}9\siunit{Н}}.
\]
\focus{Порядок действий:} перевести единицы в СИ $\rightarrow$ подобрать формулу ускорения $\rightarrow$ вычислить $F_{\text{рез}}=ma$ $\rightarrow$ проверить единицы и округление.

\heading{6. Сила, если известны две скорости и путь}
Если известно перемещение $s$ вдоль направления движения, а время отсутствует, используйте
\[
 v^2-v_0^2=2as\quad\Longrightarrow\quad
\boxed{a=\frac{v^2-v_0^2}{2s}},\quad s>0.
\]
Второй закон Ньютона даёт
\[
\boxed{F_{\text{рез}}=m\frac{|v^2-v_0^2|}{2s}}.
\]
Модуль нужен, когда спрашивают именно \emph{величину} равнодействующей силы. Если просят проекцию, знак ускорения нужно сохранить.

\heading{7. Краткое повторение: движение по окружности}
Для равномерного движения по окружности радиуса $R>0$:
\[
\nu=\frac{N}{t}=\frac1T,\qquad
\omega=\frac{2\pi}{T}=2\pi\nu,\qquad
 v=\frac{2\pi R}{T}=\omega R.
\]
\[
 a_{\text{ц}}=\frac{v^2}{R}=\omega^2R,\qquad
 F_{\text{рез}}=ma_{\text{ц}}.
\]
$T$ --- период (с), $\nu$ --- частота (Гц), $N$ --- число полных оборотов, $\omega$ --- угловая скорость (рад/с). Центростремительное ускорение и равнодействующая направлены \emph{к центру}, даже если модуль скорости не изменяется.
\clearpage

\begin{center}
{\huge\bfseries\color{darkaccent} Алгоритм: найдите равнодействующую}\par
\vspace{15mm}
\begin{tikzpicture}[
node distance=13mm and 16mm,
start/.style={ellipse,draw=accent,line width=.9pt,align=center,minimum width=38mm,minimum height=10mm,font=\normalsize},
action/.style={rectangle,rounded corners=2pt,draw=accent,line width=.9pt,align=center,text width=53mm,minimum height=12mm,font=\normalsize,inner sep=6pt},
condition/.style={diamond,aspect=1.8,draw=accent,line width=.9pt,align=center,text width=32mm,inner sep=2pt,font=\normalsize},
arrow/.style={-{Stealth[length=3mm]},line width=.95pt,draw=darkaccent},
edge label/.style={font=\small,fill=white,inner sep=2pt}
]
\node[start] (begin) {Начало};
\node[action,below=of begin] (target) {Требуется найти $\vec F_{\text{рез}}$ или $F_{\text{рез},x}$};
\node[condition,below=15mm of target] (known) {Ускорение уже известно?};
\node[action,right=26mm of known,text width=61mm] (select) {Вычислите ускорение из данных задачи:\\[2mm]
$\displaystyle a=\frac{\Delta v}{\Delta t}$\\[2mm]
$\displaystyle a=\frac{2s}{t^2}$ при $v_0=0$\\[2mm]
$\displaystyle a=\frac{v^2-v_0^2}{2s}$};
\node[action,below=39mm of known] (compute) {Найдите равнодействующую:\\ $\vec F_{\text{рез}}=m\vec a$\\ или $F_{\text{рез},x}=ma_x$};
\node[start,below=of compute] (end) {Проверьте единицы и знак};
\draw[arrow] (begin)--(target);
\draw[arrow] (target)--(known);
\draw[arrow] (known)--node[edge label,left]{да}(compute);
\draw[arrow] (known.east)--node[edge label,above]{нет}(select.west);
\draw[arrow] (select.south) |- ([yshift=8mm]compute.east) -- (compute.east);
\draw[arrow] (compute)--(end);
\end{tikzpicture}
\end{center}
\clearpage

{\huge\bfseries\color{darkaccent} Тренировочный блок}\par
\vspace{4pt}
{\bfseries Путь А · 10 заданий}\par
\vspace{2pt}
\small\color{softgray}Решайте последовательно. В ответах указывайте единицы измерения; для силы различайте модуль и проекцию.\normalsize\color{textgray}
\begin{enumerate}
\item Равнодействующая сила, действующая на тележку массой $3\siunit{кг}$, равна $12\siunit{Н}$. Найдите модуль ускорения.
\item Тело массой $5\siunit{кг}$ движется с ускорением $1{,}6\siunit{м/с\textsuperscript{2}}$. Найдите модуль равнодействующей силы.
\item Тело движется равномерно и прямолинейно. Одна сила равна $8\siunit{Н}$ и направлена вправо, а других сил, кроме второй, нет. Найдите модуль и направление второй силы.
\item Скорость лыжника массой $75\siunit{кг}$ увеличилась на $8\siunit{м/с}$ за $4\siunit{с}$ по прямой. Найдите модуль равнодействующей силы.
\item На прямолинейном участке графика скорости автомобиля массой $1500\siunit{кг}$: при $t_1=2\siunit{с}$ проекция скорости $v_{x1}=12\siunit{м/с}$, при $t_2=8\siunit{с}$ --- $v_{x2}=3\siunit{м/с}$. Найдите $F_{\text{рез},x}$ и модуль равнодействующей.
\item Действующую на тело равнодействующую увеличили в $3$ раза, а массу --- в $2$ раза. Во сколько раз изменился модуль ускорения?
\item Массу тела уменьшили в $3$ раза, а модуль ускорения увеличили в $2$ раза. Найдите отношение $F_2/F_1$ и укажите, во сколько раз изменилась равнодействующая.
\item Брусок массой $0{,}60\siunit{кг}$ начал скользить из состояния покоя и при постоянном ускорении прошёл $32\siunit{см}$ по прямой за $0{,}40\siunit{с}$. Найдите модуль равнодействующей силы.
\item Тело массой $4\siunit{кг}$ увеличило скорость с $2$ до $6\siunit{м/с}$, пройдя $8\siunit{м}$ по прямой с постоянным ускорением. Найдите модуль равнодействующей силы.
\item Тело массой $2\siunit{кг}$ равномерно движется по окружности радиуса $2\siunit{м}$ с периодом $4\siunit{с}$. Найдите модуль равнодействующей силы. Ответ выразите через $\pi$, затем округлите до десятых.
\end{enumerate}
\clearpage

{\huge\bfseries\color{darkaccent} Ответы для самопроверки}\par
\vspace{8pt}
\renewcommand{\arraystretch}{1.6}
\begin{tabularx}{\linewidth}{@{}cX@{}}
\toprule
\textbf{№} & \textbf{Ответ} \\
\midrule
1 & $4\siunit{м/с\textsuperscript{2}}$ \\
2 & $8\siunit{Н}$ \\
3 & $8\siunit{Н}$, влево \\
4 & $150\siunit{Н}$ \\
5 & $F_{\text{рез},x}=-2250\siunit{Н}$; модуль $2250\siunit{Н}$ \\
6 & Увеличился в $1{,}5$ раза \\
7 & $F_2/F_1=\frac23$; уменьшилась в $1{,}5$ раза \\
8 & $2{,}4\siunit{Н}$ \\
9 & $8\siunit{Н}$ \\
10 & $\pi^2\siunit{Н}\approx9{,}9\siunit{Н}$ \\
\bottomrule
\end{tabularx}
\vspace{12pt}
\heading{Перед следующим занятием}
Проверьте, что Вы можете без подсказок: (1) сформулировать второй закон Ньютона; (2) найти ускорение из графика скорости, времени или перемещения; (3) сравнить два состояния через отношения величин; (4) объяснить, почему при равномерном прямолинейном движении равнодействующая равна нулю.
\end{document}
